% !TeX program = xelatex
\documentclass[UTF8,fontset=fandol,10pt,a4paper]{ctexart}
\usepackage[left=22mm,right=22mm,top=23mm,bottom=22mm,headheight=14pt]{geometry}
\usepackage{fontspec}
\setmainfont{TeX Gyre Heros}
\setsansfont{TeX Gyre Heros}
\setmonofont{Latin Modern Mono}
\renewcommand{\familydefault}{\sfdefault}
\usepackage[table]{xcolor}
\usepackage{graphicx,tikz}
\usetikzlibrary{arrows.meta,positioning,fit,backgrounds,calc}
\usepackage{booktabs,tabularx,array,enumitem,float,caption,fancyhdr,hyperref}
\definecolor{blue}{HTML}{2869A9}
\definecolor{ink}{HTML}{293340}
\definecolor{muted}{HTML}{647181}
\definecolor{line}{HTML}{C5D4E2}
\definecolor{pale}{HTML}{F0F5FA}
\definecolor{warm}{HTML}{FFF7E9}
\hypersetup{colorlinks=true,urlcolor=blue,linkcolor=blue,pdftitle={Redflag RHI 与 Vulkan 接入设计},pdfauthor={Redflag}}
\color{ink}
\setlength{\parindent}{0pt}
\setlength{\parskip}{6pt}
\linespread{1.16}
\setlist[itemize]{leftmargin=1.5em,itemsep=2pt,topsep=3pt}
\setlist[enumerate]{leftmargin=1.8em,itemsep=3pt,topsep=3pt}
\renewcommand{\arraystretch}{1.38}
\newcolumntype{Y}{>{\raggedright\arraybackslash}X}
\newcolumntype{L}[1]{>{\raggedright\arraybackslash}p{#1}}
\ctexset{section={format=\Large\bfseries\color{blue},beforeskip=16pt,afterskip=8pt},subsection={format=\large\bfseries\color{blue},beforeskip=11pt,afterskip=5pt}}
\captionsetup{font=small,labelfont=bf,textfont={color=muted},skip=7pt}
\pagestyle{fancy}\fancyhf{}
\fancyhead[L]{\small\color{muted}Redflag · RHI / Vulkan · 设计方案}
\fancyhead[R]{\small\color{muted}2026-11-21}
\fancyfoot[C]{\small\color{muted}\thepage}
\renewcommand{\headrulewidth}{0pt}
\newcommand{\note}[1]{{\small\color{muted}#1\par}}
\newcommand{\code}[1]{{\small\texttt{#1}}}
\newcommand{\ruleline}{\par\vspace{4pt}{\color{line}\hrule}\vspace{8pt}}
\newcommand{\source}[2]{{\small\href{#1}{[#2]}}}
\tikzset{
 module/.style={draw=line,fill=pale,rounded corners=2pt,align=center,inner sep=7pt,text width=3.8cm,minimum height=1.15cm,font=\small},
 smallbox/.style={module,text width=2.55cm,minimum height=.95cm,font=\footnotesize},
 wide/.style={module,text width=9.5cm},
 planned/.style={module,fill=warm,dashed},
 flow/.style={-{Stealth[length=2mm]},draw=muted,line width=.65pt},
 own/.style={flow,dashed},
 label/.style={font=\scriptsize,fill=white,inner sep=2pt,text=muted}}

\usepackage{amsmath,amssymb,listings}
\hypersetup{pdftitle={Redflag 渲染公式与实现}}
\fancyhead[L]{\small\color{muted}Redflag · 渲染公式与实现}
\lstset{language=C++,basicstyle=\small\ttfamily,keywordstyle=\color{blue},
  commentstyle=\color{muted},backgroundcolor=\color{pale},frame=single,
  rulecolor=\color{line},framesep=6pt,breaklines=true,columns=fullflexible,
  keepspaces=true,showstringspaces=false,aboveskip=7pt,belowskip=7pt}
\newcommand{\loc}[2]{\par{\footnotesize\color{muted}源码：\path{#1}；\code{#2}。}\par}
\newcommand{\unit}[1]{\,\mathrm{#1}}
\begin{document}
\zihao{5}\raggedright
{\Huge\bfseries 渲染公式与实现}\par\vspace{5pt}
{\Large\color{muted}坐标、材质与间接光的计算过程}\par
\ruleline
本文对应当前 OpenGL 实现。代码摘自项目，示意图和数值例为本文编制。先说明每段计算的输入与目的，再把公式对应到语句。矩阵使用列向量，乘积从右向左作用；项目没有固定米制，位置、相机远近面和光源半径须使用同一种场景长度单位。

\section{模型位置与屏幕位置}
模型文件中的点只说明它相对于模型原点的位置。要把它画到屏幕，还需知道模型放在场景哪里、相机在哪里，以及相机如何投影。

\begin{figure}[H]\centering
\begin{tikzpicture}[node distance=10mm]
\node[smallbox] (m) {模型点 $p_m$\\相对模型原点};
\node[smallbox,right=of m] (w) {世界点 $p_w$\\整个场景的坐标};
\node[smallbox,right=of w] (v) {观察点 $p_v$\\相机前方为负 $Z$};
\draw[flow] (m)--node[label]{$M$}(w);
\draw[flow] (w)--node[label]{$V$}(v);
\node[smallbox,below=of v] (c) {裁剪坐标 $p_c$\\保留四个分量};
\node[smallbox,left=of c] (n) {标准化设备坐标\\NDC，无长度单位};
\draw[flow] (v)--node[label,right]{$P$}(c);
\draw[flow] (c)--node[label]{除以 $p_c.w$}(n);
\end{tikzpicture}
\caption{位置变换示意。视口随后把 NDC 的横纵坐标映射成像素位置。}
\end{figure}

$M$ 是模型到世界的矩阵，$V$ 是世界到相机的矩阵，$P$ 是投影矩阵。用第四分量为 1 的点表示位置，平移就能写进矩阵乘法：
\[
 p_w=M(p_m,1)^{\mathsf T},\qquad p_v=Vp_w,\qquad
 p_c=Pp_v=PVM(p_m,1)^{\mathsf T}.
\]
GLSL 中 \code{worldPos} 对应 $p_w$；\code{u\_ViewProjection} 已在 CPU 上合成 $PV$。最后一行交给 OpenGL 的仍是裁剪坐标，透视除法由后续管线完成。
\begin{lstlisting}
vec4 worldPos = model * vec4(a_Position, 1.0);
v_WorldPos = worldPos.xyz;
\end{lstlisting}
\begin{lstlisting}
gl_Position = u_ViewProjection * worldPos;
\end{lstlisting}
\loc{RedflagForge/assets/shaders/Renderer3D_Mesh.glsl}{main（顶点阶段）}
\note{依据：\source{https://learnopengl.com/Getting-started/Coordinate-Systems}{LearnOpenGL：Coordinate Systems}。模型节点的局部矩阵还会先乘到实体矩阵右侧，见下一页。}

\newpage
\subsection{变换顺序与数值核对}
Redflag 的模型可以包含多个节点。导入时累积父子节点变换，提交时再乘实体在场景中的变换；否则模型零件会失去相对位置。
\begin{lstlisting}
glm::mat4 finalTransform = transform * subMesh.LocalTransform;
\end{lstlisting}
\loc{Redflag/src/Redflag/Renderer/Renderer3D.cpp}{DrawModel}
这里 \code{LocalTransform} 是导入时积累到子网格的变换；\code{transform} 来自实体。一般有 $TS\ne ST$：先把物体缩放，再移动到场景位置，与先移动再把位移一起缩放，结果不同。

取模型点 $(0.5,1,-2)$，先沿 $X$ 放大两倍，再沿世界 $X$ 移动 3。相机在 $(0,0,3)$，朝负 $Z$，不旋转。
\[
 M=T(3,0,0)S(2,1,1),\quad
 p_w=(4,1,-2,1)^{\mathsf T},\quad p_v=(4,1,-5,1)^{\mathsf T}.
\]
再取垂直视角 $90^\circ$、宽高比 1、近面 $n=1$、远面 $f=9$。本项目采用 OpenGL 默认深度约定，投影矩阵为
\[
 P=\begin{pmatrix}
 1&0&0&0\\0&1&0&0\\
 0&0&-\frac{f+n}{f-n}&-\frac{2fn}{f-n}\\
 0&0&-1&0
 \end{pmatrix}
 =\begin{pmatrix}1&0&0&0\\0&1&0&0\\0&0&-1.25&-2.25\\0&0&-1&0\end{pmatrix}.
\]
于是 $p_c=(4,1,4,5)^{\mathsf T}$，除以 5 得到 NDC $(0.8,0.2,0.8)$。在左下为原点、原点偏移为零的 $100\times100$ 视口里：
\[
 x=\frac{0.8+1}{2}\,100=90,\quad
 y=\frac{0.2+1}{2}\,100=60,\quad d=\frac{0.8+1}{2}=0.9.
\]
这些是连续的窗口坐标；光栅化再决定哪些像素被三角形覆盖。深度 $0.9$ 不是距离 $0.9$，也不是 90\% 的相机远平面距离。

\subsection{实现前提}
模型矩阵是仿射变换，所以 \code{worldPos.w} 为 1，读取 \code{worldPos.xyz} 合理。投影矩阵不能当作模型矩阵放到这里。投影后的点还必须先通过裁剪，不能把任意 $w=0$ 或相机后方的点直接换成屏幕坐标。

\textbf{验证。}本例的矩阵乘法由随文 \path{verify.py} 复算，深度恢复部分会把同一点还原回来。项目的实体 ID 回读测试用于验证实际绘制覆盖，它不代替任意模型节点层级的完整导入验证。

\note{依据：\source{https://pbr-book.org/4ed/Geometry_and_Transformations/Applying_Transformations}{PBRT 第四版 §3.10：Applying Transformations}；投影约定同前页 LearnOpenGL。数值为本文自行设定。}

\newpage
\section{非均匀缩放后的法线与切线}
表面被拉长后，光照使用的法线仍应垂直于表面。切线沿着表面延伸，跟着几何变换；法线需要由这个垂直关系反推。

令 $A$ 为模型矩阵左上角的 $3\times3$ 线性部分。原切线 $t$ 与原法线 $n$ 满足 $n^{\mathsf T}t=0$，变换后 $t'=At$。要保持垂直，就需要
\[
 {n'}^{\mathsf T}At=0.
 \quad\text{取 }n'=A^{-\mathsf T}n,\quad
 (A^{-\mathsf T}n)^{\mathsf T}At=n^{\mathsf T}A^{-1}At=0.
\]
因此代码先对法线用逆转置，再单位化。切线用 $A$ 变换，并减去它沿法线的分量；这样两条轴重新垂直。
\begin{lstlisting}
mat3 axes = mat3(model);
vec3 N = normalize(transpose(inverse(axes)) * a_Normal);
vec3 T = axes * a_Tangent;
T = normalize(T - dot(T, N) * N);
v_Normal = N;
v_TBN = mat3(T, cross(N, T), N);
\end{lstlisting}
\loc{RedflagForge/assets/shaders/Renderer3D_Mesh.glsl}{main（顶点阶段）}

\begin{figure}[H]\centering
\begin{tikzpicture}[scale=1.05]
\draw[thick,color=muted] (-2,-1)--(2,1) node[right] {$t'\parallel(2,1)$};
\draw[flow,blue,line width=1pt] (0,0)--(.7,-1.4) node[below] {正确 $n'\parallel(0.5,-1)$};
\draw[flow,draw=orange] (0,0)--(1.4,-.7) node[right] {直接缩放 $An$};
\fill (0,0) circle(1.5pt);
\end{tikzpicture}
\caption{$X$ 放大两倍后的局部截面。橙色方向已经不垂直于表面。}
\end{figure}

\textbf{数值例。}取 $t=(1,1,0)/\sqrt2$、$n=(1,-1,0)/\sqrt2$、$A=\operatorname{diag}(2,1,1)$。错误做法的点积 $(An)\cdot(At)=1.5$；逆转置的点积为 0。单位化后 $N\approx(0.4472,-0.8944,0)$，$T\approx(0.8944,0.4472,0)$。

\textbf{前提与边界。}法线和切线是无长度单位的方向；$A$ 必须可逆，零缩放不满足条件。法线贴图中的方向经 $[T,B,N]$ 转到世界空间。当前网格只保存三分量切线，没有单独保存镜像 UV 的朝向符号；\code{cross(N,T)} 不能覆盖所有镜像 UV 情况。缺失或退化的切线也需要正确的导入约定。

\note{依据：\source{https://pbr-book.org/4ed/Geometry_and_Transformations/Applying_Transformations}{PBRT 第四版 §3.10.3：Normals}；\source{https://learnopengl.com/Lighting/Basic-Lighting}{LearnOpenGL：Basic Lighting / normal matrix}。垂直性和数值由 \path{verify.py} 检查；不把这个算例描述为所有模型法线贴图已通过视觉回归。}

\newpage
\section{直接光照的 GGX 与 Schlick 近似}
一盏灯照到表面后，只有一部分光会朝相机反射。代码分别计算照来的光，以及材质怎样分配反射光。后者用双向反射分布函数（BRDF）描述，它的理论单位是 $\mathrm{sr}^{-1}$；RGB 反射率本身没有单位。

\begin{tabularx}{\textwidth}{L{2.8cm}Y}
\toprule
符号 & 本文与代码的约定 \\
\midrule
$N,V,L$ & 世界空间单位方向：法线、表面指向相机、表面指向光源。\\
$H$ & 半程方向，$H=(V+L)/\lVert V+L\rVert$；满足该朝向的微小镜面能把 $L$ 反射到 $V$。\\
$c,m,r$ & 线性 RGB 底色、金属度、粗糙度；标量 $m,r$ 通常在 $[0,1]$。\\
$F_0$ & 正入射时的 RGB 反射率；代码按 $m$ 在 $0.04$ 和底色 $c$ 之间插值。\\
\bottomrule
\end{tabularx}

\subsection{从反射积分到一盏灯}
理论上，出射辐亮度为半球内各个入射方向贡献的积分：
\[
 L_o(V)=\int_\Omega f_r(L,V)L_i(L)\max(N\cdot L,0)\,\mathrm d\omega.
\]
$L_i,L_o$ 的物理单位为 $\mathrm{W\,m^{-2}\,sr^{-1}}$，$\mathrm d\omega$ 是立体角。项目的离散灯光把颜色乘强度放进 \code{radiance}，但没有把这些数值标定成上述物理单位。这里只能把输出理解为线性 HDR 相对亮度。

对一盏灯，当前代码使用
\[
 f_r=\underbrace{(1-F)(1-m)\frac{c}{\pi}}_{\text{漫反射}}
 +\underbrace{\frac{DGF}{4(N\cdot V)(N\cdot L)}}_{\text{镜面反射}},
 \qquad L_o^{\rm light}=f_r\,q\,\max(N\cdot L,0)\,s.
\]
$q$ 是代码中的灯光颜色与强度，$s\in[0,1]$ 是阴影可见度。背向表面的点积在代码中截为 0。离散光源的贡献已包含在 $q$ 中，不应再随手乘一个任意的立体角步长。

\subsection{三个因子的职责}
$D$ 决定微小表面朝向 $H$ 的集中程度，控制高光形状；$G$ 估算微小表面互相遮挡后还可见多少；$F$ 估算该入射角的反射比例。三者相乘后还需要分母中的方向变换因子，不能只把 $D$ 当作最终高光亮度。

\note{依据：\source{https://learnopengl.com/PBR/Theory}{LearnOpenGL：PBR Theory / reflectance equation}；\source{https://pbr-book.org/4ed/Reflection_Models/Roughness_Using_Microfacet_Theory}{PBRT 第四版 §9.6：微表面模型}。PBRT 的完整分布与遮蔽实现不等同于下一页的实时近似。}

\newpage
\subsection{公式对应到当前实现}
Redflag 使用各向同性 GGX 分布，取 $\alpha=r^2$：
\[
 D=\frac{\alpha^2}{\pi\bigl((N\cdot H)^2(\alpha^2-1)+1\bigr)^2},\quad
 F=F_0+(1-F_0)(1-H\cdot V)^5.
\]
$D$ 不是直接以立体角为测度的单位概率密度；微表面朝向分布满足 $\int D(h)(N\cdot h)\,\mathrm d\omega_h=1$。平方粗糙度是本实现的参数映射。Schlick 的五次幂是近似式，并非从当前几行代码推导出的精确 Fresnel 方程。
\begin{lstlisting}
float a  = roughness * roughness;
float a2 = a * a;
\end{lstlisting}
\begin{lstlisting}
return F0 + (1.0 - F0) * pow(clamp(1.0 - cosTheta, 0.0, 1.0), 5.0);
\end{lstlisting}
直接光的遮蔽采用 Schlick-GGX：
\[
 k=\frac{(r+1)^2}{8},\qquad
 G_1(x)=\frac{x}{x(1-k)+k},\qquad
 G=G_1(\max(N\cdot V,0))G_1(\max(N\cdot L,0)).
\]
\loc{RedflagForge/assets/shaders/Renderer3D_Mesh.glsl}{DistributionGGX / GeometrySchlickGGX / GeometrySmith / FresnelSchlick}
\begin{lstlisting}
vec3 kD = (vec3(1.0) - F) * (1.0 - metallic);
\end{lstlisting}
\begin{lstlisting}
return (kD * albedo / PI + specular) * radiance * NdotL * shadow;
\end{lstlisting}
\loc{RedflagForge/assets/shaders/Renderer3D_Mesh.glsl}{CalcDirectionalLight}

\textbf{数值例。}令 $N=V=L=(0,0,1)$、$r=0.5$、$m=0$、$c=0.5$、$q=3$、$s=1$。则 $D=16/\pi\approx5.093$，$G=1$，$F=0.04$。按代码分母额外加 $10^{-4}$，镜面项约 $0.05093$，漫反射项约 $0.15279$，单通道输出约 $0.61115$。此时尚未加入环境光、RSM 或后处理。

\textbf{工程近似。}当前粗糙度被限制为 $[0.089,1]$，GGX 分母另有下限，镜面项上限为 80；这些限制改变了高光能量，不等于严格能量守恒。点光源还使用 $(1-(d/R)^2)^2$ 的有限半径窗口，没有物理上的 $1/d^2$ 项。IBL 的 BRDF 查表使用另一种 $k=r^2/2$，不能与这里的直接光 $k$ 混用。

\note{依据：\source{https://learnopengl.com/PBR/Lighting}{LearnOpenGL：PBR Lighting}；GGX 的定义对照前页 PBRT。这里只摘项目表达式，分母下限与高光限制是本项目实现，不归于教材。数值由 \path{verify.py} 复算，未进行光度标定或完整场景的物理精度评估。}

\newpage
\section{从屏幕深度恢复表面位置}
SSAO 需要知道表面附近的采样点是否被挡住。主场景已经存下深度，就可以逆着投影过程找回表面位置，不必再给每个像素存一份三维坐标。

输入 $(u,v)$ 是纹理坐标，范围 $[0,1]$；$d$ 是深度附件中的无量纲值。Redflag 当前没有使用 Vulkan 的 $[0,1]$ 裁剪深度约定，也没有反向深度。先还原 OpenGL 的 NDC：
\[
 p_n=(2u-1,\;2v-1,\;2d-1,\;1)^{\mathsf T}.
\]
由于原来执行过透视除法，逆矩阵得到的仍是一个齐次表示；最后必须再除以其第四分量。
\[
 q=P^{-1}p_n,\qquad p_v=q.xyz/q.w.
\]
\begin{lstlisting}
vec4 ndc   = vec4(uv * 2.0 - 1.0, depth01 * 2.0 - 1.0, 1.0);
vec4 viewH = u_InvProjection * ndc;
return viewH.xyz / viewH.w;
\end{lstlisting}
\loc{RedflagForge/assets/shaders/SSAO.glsl}{ReconstructViewPos}

若需要世界位置，把 $P^{-1}$ 换成 $(PV)^{-1}$；RSM 间接光正是这样读取相机深度，避免把相机坐标与灯光世界坐标混在一起。
\begin{lstlisting}
vec4 ndc    = vec4(uv * 2.0 - 1.0, depth * 2.0 - 1.0, 1.0);
vec4 worldH = u_InvViewProj * ndc;
return worldH.xyz / worldH.w;
\end{lstlisting}
\loc{RedflagForge/assets/shaders/RSMIndirect.glsl}{ReconstructWorldPos}

\subsection{正向与逆向计算闭合}
沿用第 2 页的点，$(u,v,d)=(0.9,0.6,0.9)$。逆投影得到 $q=(0.8,0.2,-1,0.2)$，除以 $0.2$ 后是 $(4,1,-5)$，与正向计算的观察坐标一致。再还原相机平移，世界位置是 $(4,1,-2)$。

只看深度时，本例透视投影还能写成
\[
 z_v=-\frac{nf}{f-d(f-n)}.
\]
代入 $n=1,f=9,d=0.9$ 得 $z_v=-5$，而 $d=0$、$d=1$ 分别恢复到 $-1$ 和 $-9$。因此深度预览中对 $d$ 直接拉伸得到的是对比灰度，不能读作真实距离。

\textbf{使用边界。}逆投影须对应写入深度的同一相机和深度约定。SSAO 将 $d\ge0.9999$ 视为背景；屏幕外的物体没有记录。RSM 历史还原静态世界点后投向上一帧，只补偿相机变化，不包含物体自身运动或前帧深度拒绝。

\note{依据：\source{https://learnopengl.com/Getting-started/Coordinate-Systems}{LearnOpenGL：投影与透视除法}、\source{https://learnopengl.com/Advanced-Lighting/SSAO}{LearnOpenGL：SSAO}。上式由本文所列投影矩阵求逆得到；\path{verify.py} 检查位置闭合与远近面，不把它当作 SSAO 或 RSM 画质验证。}
\newpage
\section{RSM 的反射通量与面积}
方向光照亮红墙后，红墙会向地面反射一部分光。RSM（反射阴影图）先保存光源能看到的表面，再把这些表面当作小光源收集贡献。这条路径只有一次间接漫反射，没有继续计算地面把光反射给第三个表面。

\begin{figure}[H]\centering
\begin{tikzpicture}[scale=.95]
\fill[pale] (1,0) rectangle (1.2,3.3);
\fill[pale] (1,0) rectangle (8.3,-.2);
\draw[thick,blue] (1,3.3)--(1,0)--(8.3,0);
\fill[red!65!black] (1,2.3) circle(2.5pt) node[left=3pt] {$y$};
\fill[ink] (5.7,0) circle(2.5pt) node[below=3pt] {$x$};
\draw[flow] (3.1,3.3)--node[label,above] {方向光}(1.05,2.33);
\draw[flow,draw=red!65!black] (1.08,2.26)--node[label,above] {$\omega=(x-y)/r$}(5.64,.04);
\draw[flow] (1,2.3)--(2,2.3) node[right] {$n_y$};
\draw[flow] (5.7,0)--(5.7,1) node[above] {$n_x$};
\draw[flow] (5.8,.1)--node[label,right] {朝相机反射}(7.2,1.7);
\node[font=\small] at (2.7,.45) {红墙};
\node[font=\small] at (7.5,.35) {地面};
\end{tikzpicture}
\caption{本文绘制的光路示意。$y$ 是发送点，$x$ 是接收点，两端的位置和法线都使用世界空间。}
\end{figure}

\subsection{从表面面积换到光源投影面积}
设方向光正交范围宽 $w=2h$，覆盖面积 $A=w^2$。光源纹理有 $W\times H$ 个像素，每个像素对应的投影面积为 $\Delta A_\perp=A/(WH)$。这些面积使用场景长度单位的平方，不是屏幕像素面积。

令 $q$ 为方向光垂直截面的辐照度，$c_l=\max(n_y\cdot l,0)$，$l$ 从表面指向光源。发送面的实际面积 $\mathrm dA_y$ 投影后变为 $\mathrm dA_\perp=c_l\,\mathrm dA_y$。用 $C_y$ 表示线性底色，当前采用漫反射率 $\rho_y=(1-m_y)C_y$，于是
\[
 \mathrm d\Phi_y=\rho_y\,q\,c_l\,\mathrm dA_y
                =\rho_y\,q\,\mathrm dA_\perp.
\]
斜面单位实际面积收到的光较少，但一个光源像素覆盖的实际面积更大，两项抵消。因此按投影面积记录反射通量时，不再乘一次 $c_l$。背光面仍需排除。

Redflag 的 RGBA16F 附件存 $D_y=\rho_y q$，即\textbf{每单位光源投影面积的反射通量}。每纹素的通量应为 $\Phi_y=D_y\Delta A_\perp$。存密度后，可以把面积留到收集阶段乘入，避免先存很小的纹素通量、再乘回纹素数量。灯光数值未做物理标定，$q$ 在项目中仍是颜色乘强度。

\note{原算法依据：Dachsbacher、Stamminger，\source{https://users.soe.ucsc.edu/~pang/160/s13/proposal/mijallen/proposal/media/p203-dachsbacher.pdf}{Reflective Shadow Maps，I3D 2005，§3.1--3.2，印刷页 204}。原文以离散纹素通量组织计算。本文的密度定义和面积整理针对 Redflag，不把二者记成同一个量。}

\newpage
\subsection{几何因子与两个 Lambert 系数}
距离为 $r=\lVert x-y\rVert$，方向 $\omega=(x-y)/r$。发送面朝接收点的余弦为 $c_y=\max(n_y\cdot\omega,0)$，接收面朝发送点的余弦为 $c_x=\max(n_x\cdot(-\omega),0)$。不考虑两点之间遮挡时，几何因子为
\[
 G(x,y)=\frac{c_y c_x}{r^2}.
\]
原论文式（1）分子用的是未经单位化的位移向量，因而写成距离四次方。两个点积各带一个 $r$，化简后仍是 $c_yc_x/r^2$。项目先得到单位方向 \code{dirN}，原来的距离平方衰减并没有少除一次距离。

Lambert 漫反射的 BRDF 是 $\rho/\pi$。发送面出射辐亮度为 $L_y=\rho_y q c_l/\pi$，再经接收面反射，有
\begin{align*}
 L_x&=\frac{\rho_x}{\pi}\int L_y\,G(x,y)\,\mathrm dA_y\\
    &=\frac{\rho_x}{\pi^2}\int D_y\,G(x,y)\,\mathrm dA_\perp.
\end{align*}
这里暂取两点可见性为 1，这是当前 RSM 的主要近似。原论文式（1）没有显式列出这里的两个 Lambert 系数，对照时需单独核算通量定义与归一化。本节保留发送、接收两端各一个 $1/\pi$，不能再从别处重复乘除。

\subsection{圆盘采样与归一化}
接收点投到光源纹理后，在 UV 半径 $R$ 的圆盘内取 $N$ 个样本。当前正交范围是正方形，这个圆盘的世界投影面积为 $A_D=A\pi R^2=\pi(wR)^2$。均匀面积采样的积分估计为
\[
 \widehat L_x=\rho_x\,\underbrace{\frac{A\pi R^2}{\pi^2N}
      \sum_{i=0}^{N-1} D_i G_i}_{\text{RSM 输出，等于接收辐照度除以 }\pi}.
\]
主绘制再乘接收材质的 $\rho_x=(1-m_x)C_x$，$C_x$ 是接收面线性底色。$\text{Intensity}$ 是额外外观倍率，默认 1。UV 越界或没有表面的样本按零贡献处理，仍除请求的 $N$，不能通过只除有效样本数把边界重新照亮。

Vogel 圆盘使用 $r_i=\sqrt{(i+0.5)/N}$。原实现分母固定为 128，只取前 16 个时，覆盖半径约为原来的 $\sqrt{16/128}\approx0.354$。现在按实际 $N$ 生成，降低质量不会同时大幅缩小覆盖范围。等面积分布使用相同权重，不能照搬原论文另一种采样分布的径向权重。

\note{Lambert 系数依据：\source{https://pbr-book.org/4ed/Reflection_Models/Diffuse_Reflection}{PBRT 第四版 §9.2：Diffuse Reflection}。RSM 几何因子对照原论文式（1）、印刷页 204。本页圆盘估计式由上述密度与等面积采样推得。}

\newpage
\subsection{源码修改与数值例}
\textbf{修改前。}以下两句取自提交 \code{1a168c3}，分别位于 RSMDepth 的片元 \code{main} 与 RSMIndirect 的 \code{ComputeRSMIndirect}。前一句多乘入射余弦，后一句只求平均，缺少圆盘投影面积与本节约定的 Lambert 系数。
\begin{lstlisting}[title={旧实现，分属两个函数}]
o_Flux     = vec4(albedo * lightColor * NdL, 1.0);
return indirect * (intensity / float(numSamples));
\end{lstlisting}

\textbf{修改后。}发送端先读取标量或贴图中的金属度。正面输出通量密度，背面输出零。下面摘出实际表达式，省略金属贴图读取和朝向判断。
\begin{lstlisting}
vec3 diffuse = baseColor.rgb * (1.0 - metallic);
fluxDensity = diffuse * u_DirLights[0].Color.rgb * u_DirLights[0].Color.w;
o_Flux     = vec4(fluxDensity, 1.0);
\end{lstlisting}
\loc{RedflagForge/assets/shaders/RSMDepth.glsl}{main（片元阶段）}

收集阶段把平均值乘圆盘面积，并计入两个 Lambert 系数。主绘制使用相同的漫反射材质约定，纯金属不再贡献这一项。
\begin{lstlisting}
float radius = sampleRadius * lightWidth;
return indirect * (intensity * radius * radius / (PI * float(numSamples)));
\end{lstlisting}
\loc{RedflagForge/assets/shaders/RSMIndirect.glsl}{ComputeRSMIndirect}
\begin{lstlisting}
Lo += texture(u_RSMResolved, screenUV).rgb * albedo * (1.0 - metallic);
\end{lstlisting}
\loc{RedflagForge/assets/shaders/Renderer3D_Mesh.glsl}{main（片元阶段）}

\textbf{数值例。}只看单通道，令 $w=10$、$R=0.2$、$q=2$、$\rho_y=0.5$，则 $D_y=1$。假设整盘的 $G=0.25$，RSM 输出为
\[
 \frac{(10\times0.2)^2}{\pi}\,(1\times0.25)=\frac1\pi\approx0.31831.
\]
接收面 $\rho_x=0.8$ 时，输出约 $0.25465$。若几何及光源覆盖一起放大两倍，面积放大四倍，$G$ 缩为四分之一，亮度保持不变。这个人为常量算例检查系数和尺度，不代表实际场景中各点的 $G$ 相同。

\newpage
\subsection{近距离处理与验证边界}
光源像素代表一小块表面，不能在两点靠近时仍把它当成无面积的无限亮点。当前用一个光源纹素的宽度 $\delta=w/W$ 限制分母：
\[
 G_{\rm used}=\frac{c_yc_x}{\max(r^2,\delta^2)},\qquad r>0.
\]
两点重合时方向无定义，代码直接跳过。这个限制替换了固定场景单位下的 \code{DIST\_BIAS=0.05} 和几何因子上限 4。新限制随光源覆盖尺度变化，但仍不是对有限面积光源的精确积分。改变 RSM 分辨率会改变这段近距离近似，不能声称已经完全分辨率无关。

\subsection{算法仍省略的内容}
两个表面之间没有再做遮挡查询。红墙可能隔着一层挡板照亮地面，这种漏光在原 RSM 中就存在，修正面积和强度不能消除。有限采样半径也会舍弃远处反射面，因此半径改变的是积分范围，不只改变噪点多少。

这里只收集第一盏方向光直接照到的表面。没有继续传播多次反射，也没有处理间接镜面反射。$(1-m)c$ 是当前漫反射近似，没有把直接光的完整 Fresnel 方向分配搬进收集积分。法线贴图细节也未参与发送面法线。不能把纯金属不再收到漫反射，解释成已经实现了金属的间接镜面光。

历史混合只根据相机重投影，没有前帧深度拒绝和物体运动信息。它能平滑一部分噪点，不能修复遮挡缺失或错误能量。调试公式时应先检查当前帧间接光，再检查多帧历史。

\subsection{复算与实测记录}
随文 \path{verify.py} 检查余弦写法等价、投影面积抵消、圆盘系数以及统一缩放后的结果。它是 CPU 数学复算，不替代实际 Shader 测试。

\subsection{从公式回到像素}

先取一个方便手算的例子，只看红色通道。设投影边长 $w=10$，纹理坐标半径 $R=0.2$，墙面收到的光 $q=2$，反射率 $\rho_y=0.5$，于是光源图上存的是 $D_y=1$。假设整个圆盘上朝向和距离合起来给出 $G=0.25$。代入上页公式，白色表面的结果为 $(10\times0.2)^2\times 0.25/\pi\approx0.31831$。地面反射率为 $0.8$ 时，朝相机的结果约为 $0.25465$。若墙、地面和灯的覆盖范围一起放大两倍，圆盘面积变成四倍，$G$ 则降为四分之一，两者相抵。随文的 \path{verify.py} 用这组数复核系数、投影面积和尺度。

显卡测试使用另一组容易判断的输入。\path{tests/Renderer3DTests.cpp} 让 RSM 纹理处处返回 $(1,0.25,0)$，两面正对、相隔 1，投影边长为 4，纹理坐标半径为 0.2。此时红色通道应为 $0.8^2/\pi\approx0.20372$，OpenGL 实际读回与公式相差不到 0.001。只把距离改为 2，亮度降为原来的四分之一。反过来改变投影宽度和纹理坐标半径、让两者乘积保持相同，结果也相同。

最后回看 Cornell 的地面 A。最早的参考图把方向光照向绿墙，光源图里红墙的反射通量为零，调高间接光强度也不会让这堵墙染红地板。把灯转向红墙，再固定同一相机、灯光和强度 5 做算法前后对照，A 点的间接增量从 $(0.07715,0.03076,0.03076)$ 变成 $(0.48218,0.13550,0.13550)$。本节开头的展示图把强度调回 1，A 点的红色增量为 0.09644。先查灯有没有照到墙，再在固定强度下比较同一块地面，就能分清灯光方向和代码改动各自带来的变化。原始读数见 \path{output/scenes/rsm/measurements.csv}。

\subsection{仍会出现的误差}

收集阶段默认墙上的样本能照到地面。若在两点之间放一块挡板，程序仍可能把红光算过去。圆盘之外的反射面则完全收不到。靠近墙脚时，把最小距离从旧版固定的 \code{DIST\_BIAS=0.05} 与几何因子上限 4，改为随纹素宽度取值，能抑制很近的亮点。光源分辨率改变后，这一近似仍会跟着变。\source{https://doi.org/10.1145/1053427.1053460}{RSM 原论文第 3.3 节讨论了遮挡近似}

目前只有第一盏方向光参与这一次间接漫反射。墙面用几何表面的插值法线，材质的法线贴图没有进入这一步。这次计算只考虑漫反射，所以纯金属既不把漫反射光送到地面，也不接收这样的光，金属的镜面反射、多次反弹和材质随角度变化的完整 Fresnel 反射都在这次估算之外。两个独立绘制测试分别检查了受光面正照与斜照、背光面与金属，以及给金属接收面加减红色输入后的结果。

相机移动时，程序把当前点投回上一帧来平滑噪点。已有的 61 项场景检查包含 16、32、128 个圆盘样本以及切换投影、光源、场景和 RSM 设置后的历史重置。普通相机移动保留历史，上一帧裁剪范围以外的点会被拒绝。同一像素上一帧是否被其他物体挡住、物体自身移动到了哪里，仍缺少对应的深度和运动数据。测试代码见 \path{tests/DepthPrepassTests.cpp} 与 \path{tests/SceneRenderTests.cpp}，运行记录见 \path{output/validation/rsm-regression.txt}。


\note{算法边界依据：原论文 §3.3 的间接光遮挡假设，以及 §7 对非漫反射扩展的讨论。实际排错过程见《Redflag 排错记录》，经典场景的图片和采集条件另见《经典场景实测》。本文不借用旧 Mask 截帧证明新的 RSM 能量计算。}

\end{document}
